\section{External declaration}%
\label{sec:global:extern_var}
\label{sec:global:external_decls}

An external declaration declares a global symbol whose definition lives outside
the package being compiled. It is written with the keyword \token{extern} and,
declaring no definition, it has neither a body nor an initial value. The
definition is supplied at link time, by an object file passed to the compiler
alongside the package, as presented in Section~\ref{sec:global:function_body}.

\subsection{Linkage}

The keyword \token{extern} can be followed by a linkage specifier written
between parentheses. Only the \token{C} linkage is available, and it means that
the symbol is named in the object file exactly as it is written in the source
code. Without a linkage specifier, the name of the symbol is mangled following
the Ymir mangling rules, from the path of the module that declares it -- which
is the form to use for a symbol defined by another Ymir package.

\begin{lstlisting}[style=coloredverbatim]
in foo;

extern (C) fn bump(a: i32)-> i32; // linked against /bump/
extern fn inYmir(a: i32)-> i32; // linked against /_Y3foo6inYmirFi32Zi32/
\end{lstlisting}

\subsection{External variables}

An external variable is declared with \token{extern}, an optional linkage
specifier, one of the keywords \token{static} or \token{lazy}, an identifier,
the token \token{:} and finally the type of the variable. The two keywords
select the kind of variable that is expected at the other end of the link.

\begin{itemize}
\item \token{static} declares a plain memory location, of the kind a C global
  variable defines. It can be mutable.
\item \token{lazy} declares a global variable following the lazy protocol
  described in Section~\ref{sec:global:global_variables}. As any lazy variable, it
  cannot be mutable unless the mutability is that of borrowed data.
\end{itemize}

\begin{lstlisting}[style=coloredverbatimError, escapechar=@]
extern (C) static mut counter: i32; // ok, a C global variable
extern lazy shared: i32; // ok, a Ymir global variable
extern lazy dmut values: [i32]; // ok, the mutability is that of the slice data

extern (C) lazy @\hb{mut}@ x: i32; // error, a lazy variable cannot be mutable
\end{lstlisting}

The keyword \token{static} is reserved for external declarations: a global
variable declared by the package being compiled is always lazy, and declaring it
\token{static} is an error. Symmetrically, the attribute \token{@thread}, which
makes a global variable thread local, only applies to a variable the package
declares itself, as an external variable is allocated by whoever defines it.

\begin{lstlisting}[style=coloredverbatimError, escapechar=$]
$\hb{static}$ X: i32 = 12; // error, only external variables may be static

@thread
lazy Y = 34; // ok, a thread local global variable

$\hb{@thread}$
extern (C) static Z: i32; // error, an external variable cannot be thread local
\end{lstlisting}

\subsection{Linking an external declaration}

In the following example, the C file \textit{lib.c} defines a variable and a
function that the Ymir package \textit{x.yr} declares as external. The object
file produced from \textit{lib.c} is passed to the compiler at the end of the
command line, exactly as the object file of an Ymir package would be.

\begin{minipage}[t][][t]{0.35\linewidth}
\begin{lstlisting}[caption=\textit{lib.c}, style=cVerb]
int counter = 41;

int bump(int a) {
  return a + 1;
}
\end{lstlisting}
\end{minipage}\hspace{5pt}%
\begin{minipage}[t][][t]{0.6\linewidth}
%% check: skip -- /bump/ and /counter/ are defined by lib.c
\begin{lstlisting}[caption=\textit{x.yr}, style=coloredverbatim]
in x;
use std::io;

extern (C) fn bump(a: i32)-> i32;
extern (C) static mut counter: i32;

fn main() {
  counter = bump(counter);
  println("counter = ", counter);
}
\end{lstlisting}
\end{minipage}

\begin{lstlisting}[style=bashVerb, escapechar=@+]
@+\textcolor{teal!80}{alice@dev:\mtilde}@+$ gcc -c lib.c -o lib.o
@+\textcolor{teal!80}{alice@dev:\mtilde}@+$ gyc x.yr lib.o -o x.exe
@+\textcolor{teal!80}{alice@dev:\mtilde}@+$ ./x.exe
counter = 42
\end{lstlisting}

\cautionbox{ Nothing checks that the type written in an external declaration
  is the type the definition actually has. The linker matches names only, so a
  wrong type is not reported by the compiler and corrupts memory at runtime. }

\vfill%
\pagebreak
